\documentclass[12pt]{article}
\usepackage[utf8]{inputenc}
\usepackage[T1]{fontenc}
\usepackage{hyperref}
\title{Thao Language Resources for K-12 Curriculum Development}
\author{Compiled by Claude Sonnet 5.5 (claude-sonnet-5-5) for Andr\'as Kornai}
\date{}
\begin{document}
\maketitle

\section{Introduction}
Thao (ISO 639-3: ssf; Glottocode thao1240) is a critically endangered Austronesian language of Taiwan, spoken in the Sun Moon Lake area of Nantou County. Sources found in search disagree on the number of speakers. Wikipedia lists about 4 native speakers (2021), the Endangered Languages Project lists 5 fluent speakers (4 L1 and 1 L2, from a 2014 personal communication with Zeitoun), and a revival committee figure of about 10 elderly speakers was also reported. Little school-oriented material turned up in the searches for this report. No Thao textbook, graded reader or school curriculum was found. The Thao ethnic group was officially recognized as the tenth indigenous people of Taiwan on 22 September 2001.

\section{Primary Educational Resources}
No primers, textbooks or school curricula for Thao were found in search results. The Thao Dictionary below is the nearest thing to a basis for teaching materials. A review noted that its definitions are in English and that it spells the interdental fricative as c, whereas the community prefers th. Teachers would need to adapt it before classroom use.

\section{Linguistic and Reference Materials}
\begin{itemize}
\item Robert Blust, \emph{Thao Dictionary} (2003), Academia Sinica, Language and Linguistics Monograph Series, about 1100 pages with an English--Thao index: \url{https://glottolog.org/resource/reference/id/143591}
\item Academia Sinica Institute of Linguistics paper on Thao revitalization and documentation (PDF; it notes that no in-depth reference grammar has appeared): \url{https://www.ling.sinica.edu.tw/upload/researcher_manager_result/1c8effa7f30e926d8eda4eb1dc225796.pdf}
\item Wikipedia overview of the Thao language: \url{https://en.wikipedia.org/wiki/Thao_language}
\end{itemize}

\section{Religious and Community Resources}
No scripture translations or community organisation sites for Thao were found in the search results.

\section{Digital and Audio Resources}
\begin{itemize}
\item Omniglot page on Thao (Thau a lalawa), with writing and orthography notes: \url{https://www.omniglot.com/writing/thao.htm}
\item Endangered Languages Project entry for Thao: \url{https://www.endangeredlanguages.com/lang/2452}
\item PARADISEC blog post that appeared in search results (possibly related to archived Thao recordings; its content was not verified): \url{https://www.paradisec.org.au/blog/wp-json/wp/v2/posts/3583}
\end{itemize}

\section{Curriculum Development Context}
Thao has very few speakers, and the search results show no school programme. A heritage curriculum would have to be built from documentary sources, mainly the Blust dictionary and recordings, in partnership with the Thao community and Academia Sinica linguists. Nothing found indicates which materials the Council of Indigenous Peoples or local schools already use. That should be checked directly with Taiwan's Council of Indigenous Peoples and Thao community bodies, as this search did not reach them.

\end{document}